\NSPage{085}%
\begin{EESegment}{EE-TGT-P085-001}
In particular, before the curl correction the covariance identity is
\end{EESegment}
\NSFormula{NS-F-P085-001}{display}{none}
\[
  \mathcal C(W_0+\mathcal L\Sigma)=\varepsilon \mathcal T_{0,*}+\Sigma+\mathcal C(\mathcal L\Sigma).
\]
\begin{EESegment}{EE-TGT-P085-002}
This makes \NSi{NS-F-P085-002}{\mathcal L} a right inverse of
\NSi{NS-F-P085-003}{D\mathcal C(W_0)}, meaning that this operation first turns the chosen stress
change into a wave correction, and the covariance derivative then turns that correction back into
the same stress change. The quadratic remainder
\NSi{NS-F-P085-004}{\mathcal C(\mathcal L\Sigma)} stays in the stress.
The same squared partition that gives the leading covariance also puts this linear inverse
together across the slow boxes. Write \NSi{NS-F-P085-005}{W_{0,\beta}} and
\NSi{NS-F-P085-006}{\mathcal L_\beta} for the local constructions before the slow cutoff is
applied. Take a real two-vector stress \NSi{NS-F-P085-007}{\sigma(r,z,t)} that does not depend
on the angular or auxiliary variables, and whose normalized representatives
\NSi{NS-F-P085-008}{Q^{2A}\sigma} belong to \NSi{NS-F-P085-009}{\mathsf M_\alpha}. Set
\end{EESegment}
\NSFormula{NS-F-P085-010}{display}{7.35}
\begin{equation}
\begin{aligned}
  W_{0,\mathrm{phys}}^{\mathrm{as}}
    &=\sum_{\beta=(\ell,a)} Q^{-A}\eta_\beta W_{0,\beta},\\
  \mathcal L_{\mathrm{phys}}^{\mathrm{as}}\sigma
    &=\sum_{\beta=(\ell,a)} Q^{-A}\eta_\beta\mathcal L_\beta(Q^{2A}\sigma).
\end{aligned}
\tag{7.35}\label{eq:7.35}
\end{equation}
\begin{EESegment}{EE-TGT-P085-003}
Each summand uses the scale \NSi{NS-F-P085-011}{Q=2^{-\ell}} from its label
\NSi{NS-F-P085-012}{\beta=(\ell,a)}. The sums are locally finite when
\NSi{NS-F-P085-013}{q>0}. Where slow supports overlap, the assigned rectangles are disjoint,
the partitions are squared, and the physical factor is \NSi{NS-F-P085-014}{Q^{-2A}}. Together,
these facts give \NSi{NS-F-P085-015}{\mathcal B(W_{0,\mathrm{phys}}^{\mathrm{as}},
\mathcal L_{\mathrm{phys}}^{\mathrm{as}}\sigma)=\sigma}. In a chart with scale
\NSi{NS-F-P085-016}{Q}, set \NSi{NS-F-P085-017}{\Sigma=Q^{2A}\sigma} and define the normalized representatives by
\end{EESegment}
\NSFormula{NS-F-P085-018}{display}{none}
\[
  W_0^{\mathrm{as}}=Q^A W_{0,\mathrm{phys}}^{\mathrm{as}},
  \qquad
  \mathcal L^{\mathrm{as}}\Sigma
  =Q^A\mathcal L_{\mathrm{phys}}^{\mathrm{as}}(Q^{-2A}\Sigma).
\]
\begin{EESegment}{EE-TGT-P085-004}
To check whether these fields belong to the wave classes, collect all coefficients with the same
label and harmonic, as in Definition~\ref{def:6.5}. The slow cutoffs have bounded overlap, and
the common-torus bounds give \NSi{NS-F-P085-019}{W_0^{\mathrm{as}}\in \mathsf W_{1/2}} and
\NSi{NS-F-P085-020}{\mathcal L^{\mathrm{as}}\Sigma\in \mathsf W_{\alpha-1/2}}, together with
\end{EESegment}
\NSFormula{NS-F-P085-021}{display}{7.36}
\begin{equation}
  \mathcal B(W_0^{\mathrm{as}},\mathcal L^{\mathrm{as}}\Sigma)=\Sigma.
\tag{7.36}\label{eq:7.36}
\end{equation}
\begin{EESegment}{EE-TGT-P085-005}
These definitions agree wherever the charts overlap. Because the input does not depend on the
auxiliary variables, the scalar coefficients can be moved outside the covariance integral, exactly
as the identity in Proposition~\ref{prop:7.6} requires.
\end{EESegment}

\begin{EESegment}{EE-TGT-P085-006}
\subsection{Exact curls, the tails left over, and the covariance remainder.}\label{sec:7.4}
For the harmonic velocity
\NSi{NS-F-P085-022}{t_m e^{ikm\Phi}}, the condition
\NSi{NS-F-P085-023}{n_\Phi\cdot t_m=0} says that the amplitude is orthogonal to the phase
gradient, and this cancels the leading term in its divergence. Derivatives of the amplitude
\NSi{NS-F-P085-024}{t_m} still contribute. Taking the full curl of the vector potential
\NSi{NS-F-P085-025}{A_m=C_m e^{ikm\Phi}} in \eqref{eq:7.38}, rather than keeping only that
leading term, gives the corrected amplitude \NSi{NS-F-P085-026}{t_m+r_m} needed for exact
incompressibility. Differentiating the exponential gives the amplitude
\NSi{NS-F-P085-027}{t_m}, while differentiating the potential coefficient
\NSi{NS-F-P085-028}{C_m} and the cylindrical frame gives \NSi{NS-F-P085-029}{r_m}. This
subsection first estimates \NSi{NS-F-P085-030}{r_m} in Lemma~\ref{lem:7.7}, then estimates the
temporal cutoff residual in \eqref{eq:7.40}. This is the unused part of the source, multiplied by
one minus the time cutoff, plus the time derivative of that cutoff multiplied by the pulse
amplitude. The subsection finally finds the covariance of a stress increment
built from the actual velocities in Corollary~\ref{cor:7.8}.
\end{EESegment}

\begin{EESegment}{EE-TGT-P085-007}
For a normalized potential \NSi{NS-F-P085-031}{A=(A_r,A_\theta,A_z)}, define
\end{EESegment}
\NSFormula{NS-F-P085-032}{display}{7.37}
\begin{equation}
  \operatorname{curl}_* A
  =\bigl(R^{-1}\partial_\theta A_z-D_zA_\theta,
         D_zA_r-D_rA_z,
         (D_r+R^{-1})A_\theta-R^{-1}\partial_\theta A_r\bigr).
\tag{7.37}\label{eq:7.37}
\end{equation}
\begin{EESegment}{EE-TGT-P085-008}
The normalized derivatives already include the chain-rule terms that appear when the potential is
evaluated at \eqref{eq:6.3}. This means that \NSi{NS-F-P085-033}{\operatorname{curl}_* A} is the
normalized physical curl of the evaluated potential.
\end{EESegment}

\begin{EESegment}{EE-TGT-P085-009}
Fix a label, a nonzero integer \NSi{NS-F-P085-034}{m}, and a harmonic coefficient
\NSi{NS-F-P085-035}{t_m\in \mathsf W_\alpha} that does not depend on
\NSi{NS-F-P085-036}{\theta} and satisfies \NSi{NS-F-P085-037}{n_\Phi\cdot t_m=0}. The
coefficient \NSi{NS-F-P085-038}{t_m} must be supported in the prescribed slow and transverse
supports and compactly inside the pulse interval, and its extension by zero must be smooth across
all those boundaries. Choose the potential so that taking its curl and differentiating the
exponential gives the amplitude \NSPage{086}%
\NSi{NS-F-P086-001}{t_m}. The potential coefficient
\NSi{NS-F-P086-002}{C_m} and the oscillating potential
\NSi{NS-F-P086-003}{A_m} are
\end{EESegment}
\NSFormula{NS-F-P086-004}{display}{7.38}
\begin{equation}
  C_m=\frac{i\,n_\Phi\times t_m}{km|n_\Phi|^2},
  \qquad A_m=C_m e^{ikm\Phi},
  \qquad \operatorname{curl}_*A_m=(t_m+r_m)e^{ikm\Phi}.
\tag{7.38}\label{eq:7.38}
\end{equation}
\begin{EESegment}{EE-TGT-P086-001}
The last identity defines \NSi{NS-F-P086-005}{r_m} as the difference between the full curl
amplitude and the prescribed transverse amplitude. Because \NSi{NS-F-P086-006}{C_m} does not
depend on \NSi{NS-F-P086-007}{\theta}, the remainder \NSi{NS-F-P086-008}{r_m} is
\end{EESegment}
\NSFormula{NS-F-P086-009}{display}{none}
\[
  r_m=\bigl(-D_z(C_m)_\theta,
             D_z(C_m)_r-D_r(C_m)_z,
             (D_r+R^{-1})(C_m)_\theta\bigr).
\]

\begin{lemma}\label{lem:7.7}
\begin{EESegment}{EE-TGT-P086-002}
For every \NSi{NS-F-P086-010}{\alpha\in\R} and every nonzero integer
\NSi{NS-F-P086-011}{m}, let \NSi{NS-F-P086-012}{t_m\in \mathsf W_\alpha} satisfy
\NSi{NS-F-P086-013}{n_\Phi\cdot t_m=0}. Assume that it is supported in the prescribed slow and
transverse supports and compactly inside the pulse interval, and that its extension by zero is smooth
across all those boundaries. The construction in \eqref{eq:7.38} gives
\NSi{NS-F-P086-014}{C_m\in \mathsf W_{\alpha+1/2}} and
\NSi{NS-F-P086-015}{r_m\in \mathsf W_{\alpha+1/2-\kappa_s}}. The velocity
\NSi{NS-F-P086-016}{(t_m+r_m)e^{ikm\Phi}=\operatorname{curl}_*A_m} is exactly divergence-free,
both for the normalized operators and after evaluation at the phase map. The potential and velocity
extend smoothly by zero across the boundaries of their supports. For the full amplitude
\NSi{NS-F-P086-017}{a_m=t_m+r_m}, its scalar product with the phase normal satisfies
\end{EESegment}
\NSFormula{NS-F-P086-018}{display}{7.39}
\begin{equation}
  ikm\,n_\Phi\cdot a_m
  =-\bigl((D_r+R^{-1})(a_m)_r+D_z(a_m)_z\bigr),
  \qquad n_\Phi\cdot a_m\in \mathsf W_{\alpha+1/2-\kappa_s}.
\tag{7.39}\label{eq:7.39}
\end{equation}
\begin{EESegment}{EE-TGT-P086-003}
The physical potentials and pressures come from the chart versions by multiplying them by
\NSi{NS-F-P086-019}{Q^{1/2-A}} and \NSi{NS-F-P086-020}{Q^{-2A}}, respectively.
\end{EESegment}
\end{lemma}

\begin{proof}
\begin{EESegment}{EE-TGT-P086-004}
When the exponential in \NSi{NS-F-P086-021}{\operatorname{curl}_*A_m} is differentiated, the
amplitude it produces is
\end{EESegment}
\NSFormula{NS-F-P086-022}{display}{none}
\[
  ikm\,n_\Phi\times C_m
  =-\frac{n_\Phi\times(n_\Phi\times t_m)}{|n_\Phi|^2}
  =t_m-\frac{n_\Phi(n_\Phi\cdot t_m)}{|n_\Phi|^2}
  =t_m.
\]
\begin{EESegment}{EE-TGT-P086-005}
Every other term in \NSi{NS-F-P086-023}{\operatorname{curl}_*A_m} comes from differentiating the
potential coefficient \NSi{NS-F-P086-024}{C_m} or the cylindrical frame. The factor
\NSi{NS-F-P086-025}{(km)^{-1}} raises the \NSi{NS-F-P086-026}{\varepsilon} exponent by
\NSi{NS-F-P086-027}{1/2}. Applying \NSi{NS-F-P086-028}{D_r} lowers it by no more than
\NSi{NS-F-P086-029}{\kappa_s}, by \eqref{eq:6.32}. Lemma~\ref{lem:7.1} gives polynomial bounds
for the coefficients of the multiplier \NSi{NS-F-P086-030}{n_\Phi/|n_\Phi|^2}. Together, these
estimates prove the two class bounds, including all derivatives of any fixed order. The normalized
divergence is
\NSi{NS-F-P086-031}{\operatorname{div}_*a=(D_r+R^{-1})a_r+R^{-1}\partial_\theta a_\theta+D_z a_z}.
Since \NSi{NS-F-P086-032}{D_r(R^{-1})=-R^{-2}}, it follows that
\NSi{NS-F-P086-033}{(D_r+R^{-1})(R^{-1}f)=R^{-1}D_r f}. Substituting \eqref{eq:7.37} gives
\end{EESegment}
\NSFormula{NS-F-P086-034}{display}{none}
\begin{align*}
  \operatorname{div}_*\operatorname{curl}_*A
   ={}&R^{-1}D_r\partial_\theta A_z-(D_r+R^{-1})D_zA_\theta\\
     &+R^{-1}D_z\partial_\theta A_r-R^{-1}\partial_\theta D_rA_z\\
     &+D_z(D_r+R^{-1})A_\theta-R^{-1}D_z\partial_\theta A_r=0,
\end{align*}
\begin{EESegment}{EE-TGT-P086-006}
where each positive term cancels its matching negative term because the operators commute. The
chain rule keeps this identity true after evaluation at the phase map. Expanding the divergence of
\NSi{NS-F-P086-035}{a_m e^{ikm\Phi}}, whose amplitude coefficient does not depend on
\NSi{NS-F-P086-036}{\theta}, gives \eqref{eq:7.39}. Dividing by
\NSi{NS-F-P086-037}{km} and applying \NSi{NS-F-P086-038}{D_r} once gives the improvement
\NSi{NS-F-P086-039}{1/2-\kappa_s} in the exponent. The coefficient bounds and cutoffs give the
claimed support and smooth extensions.
\end{EESegment}
\end{proof}

\begin{EESegment}{EE-TGT-P086-007}
Build real velocities by pairing conjugate terms. For example,
\NSi{NS-F-P086-040}{t\cos(k\Phi)}, with real \NSi{NS-F-P086-041}{t}, has coefficients
\NSi{NS-F-P086-042}{t/2} at the two harmonics \NSi{NS-F-P086-043}{m=1,-1}. The matching
potential and pressure coefficients obey the same conjugacy condition. In particular, the pressure
that goes with a real cosine velocity mode keeps the factor \NSi{NS-F-P086-044}{i/(km)} from
\eqref{eq:7.13}.
\end{EESegment}

\begin{EESegment}{EE-TGT-P086-008}
\paragraph{The temporal cutoff and its residual.}
For the solution in Proposition~\ref{prop:7.2}, set
\NSi{NS-F-P086-045}{\widehat t_m=\psi t_m} and
\NSi{NS-F-P086-046}{\widehat\pi_m=\psi\pi_m}, then apply Lemma~\ref{lem:7.7} to
\NSi{NS-F-P086-047}{\widehat t_m}. Write \NSi{NS-F-P086-048}{\widehat r_m} for the curl
remainder of this cutoff amplitude. Before the solution is multiplied by
\NSi{NS-F-P086-049}{\psi}, it may reach both ends of the pulse interval. The cutoff makes the
potential and pressure zero on an open neighborhood of each endpoint. On the labelled rectangle
\NSPage{087}%
\NSi{NS-F-P087-001}{0\leq v\leq L_s}, the product rule in \eqref{eq:7.5} gives
\end{EESegment}
\NSFormula{NS-F-P087-002}{display}{7.40}
\begin{equation}
  \widehat t_m'+\mathcal K\widehat t_m+m^2d\widehat t_m
  +ikm n_\Phi\widehat\pi_m+f_m
  =(1-\psi)f_m+\psi't_m.
\tag{7.40}\label{eq:7.40}
\end{equation}
\begin{EESegment}{EE-TGT-P087-001}
Equation~\eqref{eq:6.16} says that
\NSi{NS-F-P087-003}{\psi=1\text{ on }|v-L_s/2|\leq L_s/5}, and both terms on the right are zero
there. Wherever either term is left, \NSi{NS-F-P087-004}{|v-L_s/2|\geq L_s/5}. Equation~\eqref{eq:7.16}
then adds a factor \NSi{NS-F-P087-005}{S_*^C e^{-cS_*}} to every estimate for a fixed coefficient
derivative. During the iteration, \NSi{NS-F-P087-006}{f_m e^{ikm\Phi}} is an actual supported
harmonic source whose extension by zero is smooth, so this local identity also describes the part
left uncancelled in the global residual.
\end{EESegment}

\begin{EESegment}{EE-TGT-P087-002}
At any fixed correction stage, changing a fixed derivative estimate to physical variables introduces
a fixed power of \NSi{NS-F-P087-007}{Q^{-1}} and a polynomial in
\NSi{NS-F-P087-008}{S_*}. Since \NSi{NS-F-P087-009}{S_*=\ell^2},
\NSi{NS-F-P087-010}{Q=2^{-\ell}}, and \NSi{NS-F-P087-011}{q\asymp Q}, every fixed choice of
\NSi{NS-F-P087-012}{M,C,N} satisfies
\end{EESegment}
\NSFormula{NS-F-P087-013}{display}{none}
\[
  q^{-N}Q^{-M}S_*^C e^{-cS_*}
  \leq C_N\exp\bigl(-c\ell^2+(M+N)\ell\log 2+2C\log\ell\bigr)
  \longrightarrow 0.
\]
\begin{EESegment}{EE-TGT-P087-003}
So every resulting physical derivative is \NSi{NS-F-P087-014}{O(q^N)} for every fixed
\NSi{NS-F-P087-015}{N}. Later corrections keep these terms as separate additive flat residuals.
The same argument handles the homogeneous cutoff tails. The pressure coefficient still satisfies
\NSi{NS-F-P087-016}{\widehat\pi_m\in \mathsf W_{\alpha+1/2}}, while the exact velocity amplitude
is \NSi{NS-F-P087-017}{\widehat t_m+\widehat r_m}, with
\NSi{NS-F-P087-018}{\widehat r_m\in \mathsf W_{\alpha+1/2-\kappa_s}}.
\end{EESegment}

\begin{EESegment}{EE-TGT-P087-004}
\paragraph{Covariance of the actual velocities.}
The linear covariance inverse was defined at the fixed primary wave. During the iteration, its
increment is added to a divergence-free wave field that already includes earlier corrections, and the
new increment has a curl remainder of its own. The exact expansion below picks out the prescribed
stress increment. It also keeps the interaction with earlier corrections, the interaction with the
new curl remainder, and the square of the new velocity.
\end{EESegment}

\begin{corollary}\label{cor:7.8}
\begin{EESegment}{EE-TGT-P087-005}
Fix \NSi{NS-F-P087-019}{\alpha,\rho\in\R}, and use the normalized representatives from
\eqref{eq:7.35}. Suppose \NSi{NS-F-P087-020}{U} is a real divergence-free wave field, including
its curl remainders, with
\end{EESegment}
\NSFormula{NS-F-P087-021}{display}{none}
\[
  U=W_0^{\mathrm{as}}+E,
  \qquad U\in \mathsf W_{1/2},
  \qquad E\in \mathsf W_\rho.
\]
\begin{EESegment}{EE-TGT-P087-006}
Let \NSi{NS-F-P087-022}{\Sigma(R,Z,T)\in \mathsf M_\alpha} be a real two-vector that does not
depend on the angular or auxiliary variables. Apply Lemma~\ref{lem:7.7} to the signed amplitudes
\NSi{NS-F-P087-023}{\mathcal L^{\mathrm{as}}\Sigma}. This gives the divergence-free increment
\NSi{NS-F-P087-024}{V=\mathcal L^{\mathrm{as}}\Sigma+R}, where
\NSi{NS-F-P087-025}{R\in \mathsf W_{\alpha-\kappa_s}} is its curl remainder. Then
\end{EESegment}
\NSFormula{NS-F-P087-026}{display}{7.41}
\begin{equation}
  \mathcal C(U+V)-\mathcal C(U)
  =\Sigma+\mathcal B(E,\mathcal L^{\mathrm{as}}\Sigma)
   +\mathcal B(U,R)+\mathcal C(V).
\tag{7.41}\label{eq:7.41}
\end{equation}
\begin{EESegment}{EE-TGT-P087-007}
The three remainder terms satisfy
\end{EESegment}
\NSFormula{NS-F-P087-027}{display}{7.42}
\begin{equation}
\begin{aligned}
  \mathcal B(E,\mathcal L^{\mathrm{as}}\Sigma)&\in \mathsf M_{\rho+\alpha-1/2},\\
  \mathcal B(U,R)&\in \mathsf M_{\alpha+1/2-\kappa_s},\\
  \mathcal C(V)&\in \mathsf M_{2\alpha-1}.
\end{aligned}
\tag{7.42}\label{eq:7.42}
\end{equation}
\begin{EESegment}{EE-TGT-P087-008}
Taking the radial divergence of any one of these remainder terms lowers the
\NSi{NS-F-P087-028}{\varepsilon} exponent by no more than another
\NSi{NS-F-P087-029}{\kappa_s}.
\end{EESegment}
\end{corollary}

\begin{proof}
\begin{EESegment}{EE-TGT-P087-009}
Expand the quadratic map and use
\NSi{NS-F-P087-030}{\mathcal B(W_0^{\mathrm{as}},\mathcal L^{\mathrm{as}}\Sigma)=\Sigma}. This gives
\end{EESegment}
\NSFormula{NS-F-P087-031}{display}{none}
\begin{align*}
  \mathcal C(U+V)-\mathcal C(U)
    &=\mathcal B(U,V)+\mathcal C(V)\\
    &=\mathcal B(W_0^{\mathrm{as}},\mathcal L^{\mathrm{as}}\Sigma)
      +\mathcal B(E,\mathcal L^{\mathrm{as}}\Sigma)+\mathcal B(U,R)+\mathcal C(V)\\
    &=\Sigma+\mathcal B(E,\mathcal L^{\mathrm{as}}\Sigma)
      +\mathcal B(U,R)+\mathcal C(V).
\end{align*}
\begin{EESegment}{EE-TGT-P087-010}
The bounds for \NSi{NS-F-P087-032}{\mathcal B(E,\mathcal L^{\mathrm{as}}\Sigma)} and
\NSi{NS-F-P087-033}{\mathcal B(U,R)} come from multiplying the stated classes. Since
\NSi{NS-F-P087-034}{\kappa_s<1/2}, it follows that
\NSi{NS-F-P087-035}{V\in \mathsf W_{\alpha-1/2}}, and so
\NSi{NS-F-P087-036}{\mathcal C(V)\in \mathsf M_{2\alpha-1}}. This proves \eqref{eq:7.41}
and \eqref{eq:7.42} without placing any positivity condition \NSPage{088}%
on an accumulated stress. The claim
about radial divergence follows from \eqref{eq:6.32}; the cylindrical terms of order zero cause no
further loss in the exponent.
\end{EESegment}
\end{proof}

\begin{EESegment}{EE-TGT-P088-001}
\section{Compactly supported corrections to the mean}\label{sec:8}
The wave construction in Section~\ref{sec:7} gives pulses whose averaged covariance produces the
leading stress target in \eqref{eq:7.26}. It also gives an inverse for the principal equation of every
nonzero angular harmonic in Proposition~\ref{prop:7.2}. The residual's angular mean, meaning its
average over the physical angle, still needs to be corrected. This mean contains quadratic products
of the waves, and it can depend on the auxiliary torus even though it does not depend on the physical
angle, as Proposition~\ref{prop:8.1} shows. This section builds the pressure, tangential stresses, and
divergence-free velocity increments needed to correct it. All of them are supported inside the active
annulus; see Proposition~\ref{prop:8.3} and Corollary~\ref{cor:8.5}.
\end{EESegment}

\begin{EESegment}{EE-TGT-P088-002}
There are two different inversion problems. For an auxiliary-averaged source, a compactly supported
solution \NSi{NS-F-P088-001}{\varphi} of
\NSi{NS-F-P088-002}{(D_r+e/R)\varphi=f}, with
\NSi{NS-F-P088-003}{e\in\{0,1,2\}}, can exist only when
\NSi{NS-F-P088-004}{\int R^e f\,dR=0}; see Lemma~\ref{lem:8.2}. For a source that depends on the
auxiliary torus, the fast time derivative can be inverted only when the source has zero auxiliary
average at each slow point, as Lemma~\ref{lem:8.6} explains. The section first derives the mean
momentum equations in Proposition~\ref{prop:8.1} and builds the radial inverse in
Lemma~\ref{lem:8.2}, including the remainder that appears when its source depends on the torus. It
then uses that inverse for pressure and divergence-free axial increments in
Proposition~\ref{prop:8.3}, and for tangential stresses in
Corollary~\ref{cor:8.5}. The integrated equations in \eqref{eq:8.16} pick out three scalar defects,
\NSi{NS-F-P088-005}{\mathcal P,\mathcal J_\theta,\mathcal J_z}, defined in Equations~\eqref{eq:8.12}
and \eqref{eq:8.15}. A final linear map chooses five velocity coefficients so that their linear
changes cancel these three defects while two velocity moments stay unchanged; this is
Lemma~\ref{lem:8.7}. Lemma~\ref{lem:8.8} calculates the nonlinear changes explicitly. The correction
cycle in Proposition~\ref{prop:9.6} uses all these constructions after every wave update.
\end{EESegment}

\begin{EESegment}{EE-TGT-P088-003}
Any radial integral without stated limits runs from zero to infinity. The averaging and projection
conventions are the ones in Equations~\eqref{eq:3.7} and \eqref{eq:3.9}. Recall that
\NSi{NS-F-P088-006}{\langle\,\cdot\,\rangle_\theta} takes the angular average of each cylindrical
component, while \NSi{NS-F-P088-007}{\langle\,\cdot\,\rangle_Y} takes the Haar average on the
auxiliary torus from Lemma~\ref{lem:6.2}, and
\NSi{NS-F-P088-008}{f^\circ=f-\langle f\rangle_Y}. So a field can be angularly invariant and still
depend on that torus. The coefficient classes \NSi{NS-F-P088-009}{\mathsf M_\alpha} and
\NSi{NS-F-P088-010}{\mathsf S_\alpha} are the classes from Definition~\ref{def:6.4}; their estimates
hold at every fixed derivative order. Unless physical variables are written out, the variables and
operators are the normalized chart versions from \eqref{eq:6.6}.
\end{EESegment}

\begin{EESegment}{EE-TGT-P088-004}
\subsection{Momentum equations in divergence form and integral constraints.}\label{sec:8.1}
Here, divergence form means writing the transport terms as radial and axial derivatives of
momentum fluxes, with the extra terms required by cylindrical coordinates kept beside them.
This subsection derives the mean momentum equations in divergence form in
Proposition~\ref{prop:8.1}, and states the two velocity moments in \eqref{eq:8.2} that the corrections
must leave unchanged. Radial integration of these equations shows which compactly supported
corrections are possible, as Proposition~\ref{prop:8.4} and Corollary~\ref{cor:8.5} explain. Writing
the equations in divergence form also shows the boundary contributions explicitly. The calculation
keeps the quadratic products of the actual divergence-free field
\NSi{NS-F-P088-011}{w}, including its curl remainders, as shown in \eqref{eq:8.1}. In a common
chart, using the normalized operators from \eqref{eq:6.6}, the actual velocity splits as
\end{EESegment}
\NSFormula{NS-F-P088-012}{display}{8.1}
\begin{equation}
  u=(b+\beta,V+v,G+\gamma)+w,
  \qquad \langle w\rangle_\theta=0,
  \qquad W_{ab}=\langle w_aw_b\rangle_\theta.
\tag{8.1}\label{eq:8.1}
\end{equation}
\begin{EESegment}{EE-TGT-P088-005}
Here \NSi{NS-F-P088-013}{(b,V,G)} is the slow base,
\NSi{NS-F-P088-014}{(\beta,v,\gamma)} is its correction that does not depend on the
physical angle, and \NSi{NS-F-P088-015}{w} is the sum of the divergence-free fields made by taking
curls of the wave potentials. The indices in \NSi{NS-F-P088-016}{W_{ab}} run over
\NSi{NS-F-P088-017}{a,b\in\{r,\theta,z\}}, so \NSi{NS-F-P088-018}{W} includes every curl
remainder. The correction components and \NSi{NS-F-P088-019}{W} extend smoothly by zero outside
the active shell. The mean correction is divergence-free, which means
\NSi{NS-F-P088-020}{(D_r+R^{-1})\beta+D_z\gamma=0}, and it must satisfy the two integral
constraints
\end{EESegment}
\NSFormula{NS-F-P088-021}{display}{8.2}
\begin{equation}
  M_\theta:=\int_0^\infty R^2\langle v\rangle_Y\,dR=0,
  \qquad
  M_z:=\int_0^\infty R\langle\gamma\rangle_Y\,dR=0.
\tag{8.2}\label{eq:8.2}
\end{equation}
\NSPage{089}%
\begin{EESegment}{EE-TGT-P089-001}
These two identities hold at every \NSi{NS-F-P089-001}{(Z,T)}. Up to normalization, they say
that the auxiliary-averaged correction has zero angular-momentum moment and zero axial-flux
integral. The base is handled separately. Compactly supported azimuthal potentials leave the
axial-flux constraint unchanged, and the finite-dimensional moment correction leaves both
constraints unchanged.
\end{EESegment}

\begin{EESegment}{EE-TGT-P089-002}
Let \NSi{NS-F-P089-002}{p_m} be the perturbation's mean pressure. Put
\NSi{NS-F-P089-003}{\Delta_0=D_r^2+R^{-1}D_r+D_z^2} and
\NSi{NS-F-P089-004}{\Sigma_a=Q^{2A}\mathcal T_{\mathrm{phys},a}} for
\NSi{NS-F-P089-005}{a=\theta,z}, where the full residual stress tensor is the one in
\eqref{eq:5.41}.
\end{EESegment}

\begin{proposition}\label{prop:8.1}
\begin{EESegment}{EE-TGT-P089-003}
Let \NSi{NS-F-P089-006}{u} have the divergence-free decomposition in \eqref{eq:8.1}, with slow
base \NSi{NS-F-P089-007}{(b,V,G)}, correction
\NSi{NS-F-P089-008}{(\beta,v,\gamma)} independent of the physical angle, and
\NSi{NS-F-P089-009}{\langle w\rangle_\theta=0}. Let \NSi{NS-F-P089-010}{p_m} be its pressure
correction independent of the physical angle, and let
\NSi{NS-F-P089-011}{\Sigma_a=Q^{2A}\mathcal T_{\mathrm{phys},a}} be the normalized components of
the base stress from \eqref{eq:5.41}. Keep the base's flat residual as a separate additive error. All
its derivatives of every fixed order vanish faster than every power of
\NSi{NS-F-P089-012}{q}. Then the angular mean of the rest of the normalized residual is
\NSi{NS-F-P089-013}{(D_rp_m-g_r,E_\theta,E_z)}, where
\end{EESegment}
\NSFormula{NS-F-P089-014}{display}{8.3}
\begin{equation}
\begin{aligned}
E_\theta={}&\mathfrak t_*v+(D_r+2/R)(bv+\beta V+\beta v+W_{r\theta})\\
 &+D_z(Gv+V\gamma+\gamma v+W_{z\theta})
   -\varepsilon(\Delta_0-R^{-2})v-(D_r+2/R)\Sigma_\theta,\\
E_z={}&\mathfrak t_*\gamma+(D_r+1/R)(b\gamma+\beta G+\beta\gamma+W_{rz})\\
 &+D_z(2G\gamma+\gamma^2+W_{zz}+p_m)-\varepsilon\Delta_0\gamma-(D_r+1/R)\Sigma_z,\\
g_r={}&-\mathfrak t_*\beta-(D_r+1/R)(2b\beta+\beta^2+W_{rr})\\
 &-D_z(b\gamma+G\beta+\beta\gamma+W_{zr})
   +(2Vv+v^2+W_{\theta\theta})/R+\varepsilon(\Delta_0-R^{-2})\beta.
\end{aligned}
\tag{8.3}\label{eq:8.3}
\end{equation}
\end{proposition}

\begin{proof}
\begin{EESegment}{EE-TGT-P089-004}
Incompressibility lets the transport terms be written in divergence form. For the cylindrical
components of any divergence-free field, the three transport expressions are
\end{EESegment}
\NSFormula{NS-F-P089-015}{display}{none}
\begin{align*}
  (u\cdot\nabla u)_r
    &=(D_r+1/R)u_r^2+R^{-1}\partial_\theta(u_\theta u_r)+D_z(u_z u_r)-u_\theta^2/R,\\
  (u\cdot\nabla u)_\theta
    &=(D_r+2/R)(u_r u_\theta)+R^{-1}\partial_\theta u_\theta^2+D_z(u_z u_\theta),\\
  (u\cdot\nabla u)_z
    &=(D_r+1/R)(u_r u_z)+R^{-1}\partial_\theta(u_\theta u_z)+D_z u_z^2.
\end{align*}
\begin{EESegment}{EE-TGT-P089-005}
The turning of the cylindrical frame produces the centrifugal term in the first row and the extra
\NSi{NS-F-P089-016}{u_ru_\theta/R} in the second. Any angular derivative integrates to zero, and
so does any product with exactly one factor of \NSi{NS-F-P089-017}{w}. Once the pure base terms
are subtracted, the fluxes left are the ones in \eqref{eq:8.3}. The angular-derivative terms in the
vector Laplacian also have zero angular average. This leaves
\NSi{NS-F-P089-018}{\Delta_0-R^{-2}} in the radial and angular components and
\NSi{NS-F-P089-019}{\Delta_0} in the axial component. By \eqref{eq:5.41}, the residual stress tensor
contributes the two radial divergences displayed in \eqref{eq:8.3}. The pulse tails left out have
nonzero angular harmonics, so their angular mean is zero. The calculation keeps the full covariance
of the divergence-free wave field, including all its curl remainders.
\end{EESegment}
\end{proof}

\begin{EESegment}{EE-TGT-P089-006}
\subsection{A compactly supported primitive for the cylindrical weights.}\label{sec:8.2}
This subsection builds a radial primitive by integrating in the radial direction along the matching
torus path, and keeps its support inside the active shell; see Lemma~\ref{lem:8.2}. A direct radial
integral may still be nonzero after the source support has ended. To stop that from happening,
subtract its full radial integral multiplied by a fixed cutoff, and keep the error created by this
subtraction. After evaluation at the phase, the radial derivative also differentiates the auxiliary
variable in \eqref{eq:6.4}, so both integrals must follow the matching path on the torus. That path is
included in the physical definition below.
\end{EESegment}

\begin{EESegment}{EE-TGT-P089-007}
Start with the radial operations in physical variables, keeping
\NSi{NS-F-P089-020}{(z,t)} and the independent angular variable fixed. Choose a slow cutoff
\NSi{NS-F-P089-021}{\chi_m(X)} that is zero on an inner collar, equals one on an outer collar, and
changes only in a fixed interior part of the active shell. Let
\NSi{NS-F-P089-022}{g(r,z,t,Y)} be a smooth source that is periodic in
\NSi{NS-F-P089-023}{Y} and extends smoothly by zero outside that shell. Set
\NSPage{090}%
\end{EESegment}
\NSFormula{NS-F-P090-001}{display}{8.4}
\begin{equation}
\begin{aligned}
  \mathcal I g(r,Y)
    &=\int_0^r g\bigl(r',Y+((r')^{d_r}-r^{d_r})v_r\bigr)\,dr',\\
  \mathcal J g(r,Y)
    &=\int_0^\infty g\bigl(r',Y+((r')^{d_r}-r^{d_r})v_r\bigr)\,dr',\\
  \mathcal I_c g&=\mathcal I g-\chi_m\mathcal J g.
\end{aligned}
\tag{8.4}\label{eq:8.4}
\end{equation}
\begin{EESegment}{EE-TGT-P090-001}
The shift in \NSi{NS-F-P090-002}{Y} follows the radial phase map while the physical variables
\NSi{NS-F-P090-003}{(z,t)} stay fixed. Below the shell,
\NSi{NS-F-P090-004}{\mathcal I g=0}; above it,
\NSi{NS-F-P090-005}{\mathcal I g=\mathcal J g}. It follows that
\NSi{NS-F-P090-006}{\mathcal I_c g} is zero when
\NSi{NS-F-P090-007}{X\leq X_a} and when \NSi{NS-F-P090-008}{X\geq X_b}, and it extends
smoothly by zero across both radial edges. The subtraction creates an error supported where
\NSi{NS-F-P090-009}{\chi_m} changes, and that error is kept in the radial equation.
\end{EESegment}

\begin{EESegment}{EE-TGT-P090-002}
For the normalized operators, use a common torus \NSi{NS-F-P090-010}{y=Y_{i_0}} and set
\NSi{NS-F-P090-011}{M=\Lambda_g^{i_0}Q^{d_r}/2}. At fixed
\NSi{NS-F-P090-012}{(Z,T)}, the same symbols mean
\end{EESegment}
\NSFormula{NS-F-P090-013}{display}{8.5}
\begin{equation}
\begin{aligned}
  \mathcal I g(R,y)
    &=\int_0^R g\bigl(R',y+M((R')^{d_r}-R^{d_r})v_r\bigr)\,dR',\\
  \mathcal J g(R,y)
    &=\int_0^\infty g\bigl(R',y+M((R')^{d_r}-R^{d_r})v_r\bigr)\,dR',\\
  \mathcal I_c g&=\mathcal I g-\chi_m\mathcal J g.
\end{aligned}
\tag{8.5}\label{eq:8.5}
\end{equation}
\begin{EESegment}{EE-TGT-P090-003}
Here \NSi{NS-F-P090-014}{\chi_m=\chi_m(QR^2/(2q))} is how the fixed physical cutoff appears in
the chart. For \NSi{NS-F-P090-015}{e\in\{0,1,2\}}, define
\end{EESegment}
\NSFormula{NS-F-P090-016}{display}{8.6}
\begin{equation}
  \mathcal D_e=D_r+e/R,
  \qquad \mathsf T_e f=R^{-e}\mathcal I_c(R^e f),
  \qquad \mathcal A_e f=R^{-e}(\partial_R\chi_m)\mathcal J(R^e f).
\tag{8.6}\label{eq:8.6}
\end{equation}
\begin{EESegment}{EE-TGT-P090-004}
At matching physical and chart points, where
\NSi{NS-F-P090-017}{r=\sqrt Q\,R}, the normalization is
\end{EESegment}
\NSFormula{NS-F-P090-018}{display}{none}
\[
  f_*=Q^a f_{\mathrm{phys}},
  \qquad
  Q^{a-1/2}r^{-e}\mathcal I_{c,\mathrm{phys}}(r^e f_{\mathrm{phys}})
   =R^{-e}\mathcal I_c(R^e f_*)=\mathsf T_e f_*.
\]
\begin{EESegment}{EE-TGT-P090-005}
Here \NSi{NS-F-P090-019}{\mathcal I_{c,\mathrm{phys}}} is the physical operator defined above.
The term \NSi{NS-F-P090-020}{\mathcal A_e f} is the cutoff remainder created by making the
primitive compactly supported.
\end{EESegment}

\begin{EESegment}{EE-TGT-P090-006}
The inverse must keep the source's decay order. The next lemma assumes that the source is multiplied
by the stated power of the radial variable, averaged over the auxiliary torus, and then integrated
over the full radial line. At every fixed value of the other two slow variables, the result must be
zero. Under that exact condition, the cutoff remainder must also be flat. The lemma gives both
estimates.
\end{EESegment}

\begin{lemma}\label{lem:8.2}
\begin{EESegment}{EE-TGT-P090-007}
Fix a correction stage \NSi{NS-F-P090-021}{j}, an exponent
\NSi{NS-F-P090-022}{\alpha}, a band, its common torus
\NSi{NS-F-P090-023}{y=Y_{i_0}}, and \NSi{NS-F-P090-024}{e\in\{0,1,2\}}. Let
\NSi{NS-F-P090-025}{f\in \mathsf M_\alpha} be a smooth scalar coefficient. In particular, for
every fixed multi-index \NSi{NS-F-P090-026}{I} of derivatives in
\NSi{NS-F-P090-027}{(R,Z,T,y)}, its smooth extension by zero satisfies
\end{EESegment}
\NSFormula{NS-F-P090-028}{display}{none}
\[
  |\partial^I f|\leq C_{j,I}\varepsilon^\alpha S_*^{B_{j,I}}\zeta\delta^{-B_{j,I}}.
\]
\begin{EESegment}{EE-TGT-P090-008}
Then \NSi{NS-F-P090-029}{\mathsf T_e f} is supported in the same active radial shell and in the
projection of the input support to \NSi{NS-F-P090-030}{(Z,T)}. Under the primitive normalization
defined above, it belongs to \NSi{NS-F-P090-031}{\mathsf M_\alpha}. The exact identity is
\end{EESegment}
\NSFormula{NS-F-P090-032}{display}{8.7}
\begin{equation}
  \mathcal D_e\mathsf T_e f=f-\mathcal A_e f.
\tag{8.7}\label{eq:8.7}
\end{equation}
\begin{EESegment}{EE-TGT-P090-009}
If \NSi{NS-F-P090-033}{\int R^e\langle f\rangle_Y\,dR=0} at every
\NSi{NS-F-P090-034}{(Z,T)}, then
\NSi{NS-F-P090-035}{\langle\mathcal A_e f\rangle_Y=0}. Also, for every fixed
\NSi{NS-F-P090-036}{I} and every integer \NSi{NS-F-P090-037}{p\geq1},
\end{EESegment}
\NSFormula{NS-F-P090-038}{display}{8.8}
\begin{equation}
  |\partial^I\mathcal A_e f|
  \leq C_{j,I,p}\varepsilon^{\alpha+p\kappa_s}S_*^{B_{j,I,p}}
         \zeta\delta^{-B_{j,I,p}}.
\tag{8.8}\label{eq:8.8}
\end{equation}
\begin{EESegment}{EE-TGT-P090-010}
For each \NSi{NS-F-P090-039}{(I,p)}, only finitely many derivatives of the input are used. The
constants and powers may depend on those indices, the stage, and the fixed profile, but not on the
band, label, or point. If \NSi{NS-F-P090-040}{f} does not depend on
\NSi{NS-F-P090-041}{Y} and \NSi{NS-F-P090-042}{\int R^e f\,dR=0} at every slow point, the cutoff
remainder is identically zero.
\end{EESegment}
